\documentclass[11pt,a4paper]{article}
\usepackage[T1]{fontenc}
\usepackage{lmodern,amsmath,amssymb,amsthm,mathtools}
\usepackage[margin=26mm]{geometry}
\usepackage{array,longtable,booktabs,xcolor,microtype}
\usepackage[colorlinks=true,linkcolor=blue!45!black,citecolor=blue!45!black,urlcolor=blue!45!black]{hyperref}
\newtheorem{theorem}{Theorem}[section]
\newtheorem{lemma}[theorem]{Lemma}
\newtheorem{proposition}[theorem]{Proposition}
\newtheorem{corollary}[theorem]{Corollary}
\theoremstyle{definition}
\newtheorem{definition}[theorem]{Definition}
\theoremstyle{remark}
\newtheorem{remark}[theorem]{Remark}
\newcommand{\Q}{\mathbb Q}
\newcommand{\Z}{\mathbb Z}
\newcommand{\Sha}{\operatorname{Sha}}
\setlength{\emergencystretch}{2em}
\setcounter{tocdepth}{1}
\title{Eliminating finite algebraic choices with one coordinate}
\author{}
\date{7 November 2026}
\begin{document}
\maketitle
\begin{abstract}
We give an explicit polynomial which replaces all finite
algebraic choices in a triangular monic presentation by one rational
coordinate. The proof includes degree drops, nonzero guards and
inactive nonreduced fibers. We extend the construction to every
number field using explicit anisotropic forms of degree $e$ in
$e^2$ variables. The results are unconditional algebraic statements;
they require no arithmetic hypothesis about definability of the integers.
\end{abstract}
\section{Finite choices and rational witnesses}
An existential polynomial construction often contains coordinates
which, once its other coordinates are fixed, have only finitely many
possible values. We give an explicit elimination of all such finite
choices using one further rational coordinate. The construction
handles nonzero conditions, drops in degree and inactive nonreduced
fibers. All products and determinants in the formula have explicit
finite index sets.

The result is algebraic and does not depend on any assertion about
Diophantine definitions of the integers. Its extension to arbitrary
number fields uses an anisotropic form to combine its residuals.
\section{An explicit polynomial for a guarded finite algebra}
\label{sec:explicit-finite-algebra}

Let $R=\mathbb Z[x_1,\ldots,x_b]$. Consider a triangular monic presentation
\[
 A=R[z_1,\ldots,z_m]/(f_1,\ldots,f_m),\qquad
 f_i\in R[z_1,\ldots,z_i],
\]
where $f_i$ is monic of degree $e_i\ge1$ in $z_i$. Its standard
monomials $\prod_i z_i^{a_i}$, $0\le a_i<e_i$, form an $R$-basis.
In particular $A$ is free of rank $n=\prod_i e_i$. For $a\in A$,
write $M_a$ for its multiplication matrix in this ordered monomial basis.
The matrices $M_a$ commute, and their entries are integer polynomials in
the base coordinates. They are completely specified by repeated reduction
using the monic equations; they need not be materialized to specify a
determinant involving them.

Fix elements $G,E_1,\ldots,E_r\in A$. We assume that in every geometric
fiber, every point at which $G$ is nonzero has local length one. This
condition holds, for example, if $G$ contains the product of the diagonal
derivatives $\partial f_i/\partial z_i$ as a factor. We seek the rational
base points admitting a rational point of $A$ with
\[
 G\ne0,\qquad E_1=\cdots=E_r=0.
\]

For $0\le\ell\le m n(n+1)/2$, define
\[
 \lambda_\ell=1+\sum_{i=1}^m\ell^i z_i,\qquad
 M=M_{G\lambda_\ell},\qquad
 H(V)=\det(I-VM)=\sum_{k=0}^n a_kV^k.
\]
Here and below $a_k\in R$ depend on the fixed integer $\ell$; always
$a_0=1$. Also put
\[
 R_i(V)=\operatorname{Tr}\bigl(M_{E_i}\operatorname{adj}(I-VM)\bigr).
\]
This is a polynomial in $R[V]$ of degree at most $n-1$.

For each $1\le d\le n$, set
\begin{align*}
 H_d(V)&=\sum_{k=0}^d a_kV^k,\\
 g&=a_d\operatorname{Res}_{V}(H_d,H'_d),\qquad L=a_d^2g,\\
 \Psi(U)&=1+\sum_{j=1}^d a_{d-j}a_d^{2j-1}g^jU^j,\\
 \widetilde R_i(U)&=(LU)^{n-1}R_i\bigl(1/(LU)\bigr),\\
 Q_{\ell,d}(\mathbf x,U)&=
 \Psi(U)^2+\sum_{k=d+1}^n a_k^2+\sum_{i=1}^r\widetilde R_i(U)^2.
\end{align*}
The resultant is the Sylvester determinant with the specified degree
bounds $d$ and $d-1$. The expression for $\widetilde R_i$ denotes
polynomial coefficient reversal, hence defines a polynomial also at
$LU=0$. Every displayed expression has integer coefficients.

\begin{theorem}\label{thm:explicit-guarded-finite-algebra}
The explicitly specified polynomial
\[
 P(\mathbf x,U)=
 \prod_{\ell=0}^{m n(n+1)/2}\ \prod_{d=1}^{n} Q_{\ell,d}(\mathbf x,U)
\]
has the following property. For every $\mathbf a\in\mathbb Q^b$,
there exists $U\in\mathbb Q$ with $P(\mathbf a,U)=0$ if and only if
there is a rational point of the fiber algebra $A_{\mathbf a}$ satisfying
$G\ne0$ and all $E_i=0$. The formula uses only one coordinate in addition
to the $b$ base coordinates; its formal algebra generators are eliminated
by the determinants and are not quantified coordinates.
\end{theorem}
\begin{proof}
Fix a rational base point. First consider one factor $Q_{\ell,d}$.
At $a_dg=0$ or $U=0$, its reciprocal residual $\Psi$ equals one.
Otherwise, with $V=1/(LU)$,
\[
 \Psi(U)=(LU)^d\frac{H_d(V)}{a_d}.
\]
Thus a rational zero of $Q_{\ell,d}$ implies $a_dgU\ne0$, all
coefficients $a_k$ with $k>d$ vanish, and $H(V)=H_d(V)=0$.
The resultant condition makes this a simple root of $H$. Since
$H(0)=1$, the root $V$ is nonzero.

The rational matrix $I-VM$ has a one-dimensional kernel. Every
multiplication matrix $M_{z_i}$ preserves it and therefore acts on
it by a rational scalar $\zeta_i$. Applying the algebra relations
to this common eigenline shows that $(\zeta_1,\ldots,\zeta_m)$ is
a rational point of the fiber. Its $G\lambda_\ell$-value is $1/V$,
so its $G$-value is nonzero.

To verify the closed conditions, write
$J=\operatorname{Tr}(\operatorname{adj}(I-VM))$. At a root of $H$,
the adjugate identity gives
$M\operatorname{adj}(I-VM)=V^{-1}\operatorname{adj}(I-VM)$.
Differentiating the determinant then yields
\[
 J=-V H'(V)\ne0.
\]
The columns of the adjugate lie in the common eigenline, so
\[
 R_i(V)=E_i(\boldsymbol\zeta)J.
\]
Since $LU\ne0$, the equation $\widetilde R_i(U)=0$ is equivalent
to $R_i(V)=0$, and hence $E_i(\boldsymbol\zeta)=0$. This proves
the reverse implication without assuming that the rest of the fiber
algebra is reduced.

Conversely, suppose a suitable rational point exists. Over an
algebraic closure, multiplication by $G\lambda_\ell$ has the
value $G(p)\lambda_\ell(p)$ at every geometric point $p$, with
multiplicity equal to its local length. If $G(p)=0$, its factor
in $\det(I-VM)$ is one, even if the local algebra is nonreduced.
If $G(p)\ne0$, its local length is one by hypothesis.

At any such point $p$, the polynomial
$G(p)(1+\sum_{i=1}^m\ell^i z_i(p))$ is nonzero as a polynomial
in $\ell$, because its constant coefficient is $G(p)\ne0$.
For distinct such points $p,q$, their difference is also nonzero:
if all coefficients vanished, first $G(p)=G(q)\ne0$, and then
all their coordinates would agree. Each of these polynomials has
degree at most $m$. With at most $n$ geometric points, there are
at most $m n(n+1)/2$ bad values of $\ell$ in total, counting zero
values and collisions. Thus one of the specified integers makes
all nonzero guarded eigenvalues distinct and nonzero.

For that $\ell$, the polynomial $H$ is squarefree of some degree
$1\le d\le n$. The given rational point provides a nonzero
rational root $V$ of $H$. Its closed residuals vanish, so the trace
identities give $R_i(V)=0$. The leading coefficient and resultant
are nonzero. Taking $U=1/(LV)$ therefore gives a zero of the
corresponding factor $Q_{\ell,d}$, and hence of $P$.
\end{proof}

The degree cases are part of the literal finite formula. They include
every exceptional drop in degree of $H$. The discriminant condition
excludes bad separating choices, and the finite bound on $\ell$ proves
that this exclusion loses no valid rational point. In particular the
construction does not leave a generic primitive-element argument or an
unperformed exceptional-stratum decomposition implicit.


\section{Extension to arbitrary number fields}
\begin{definition}
A homogeneous polynomial $A\in F[X_1,\ldots,X_s]$ is anisotropic
over $F$ if its only zero in $F^s$ is the zero tuple.
\end{definition}

\begin{lemma}\label{fa:anisotropic}
For every number field $F$ and every integer $e\ge1$ there is an
explicit anisotropic form of degree $e$ in $e^2$ variables over $F$.
\end{lemma}
\begin{proof}
Choose a nonzero prime ideal $\mathfrak p$ of $\mathcal O_F$, an
element $\pi\in\mathfrak p\setminus\mathfrak p^2$, and a monic
polynomial $g\in\mathcal O_F[Z]$ of degree $e$ whose reduction is
irreducible over the residue field. Such a polynomial is obtained
by lifting the coefficients of any irreducible degree-$e$ polynomial
over that finite field. Let $N_g(\mathbf X)$ be the determinant of
multiplication by $\sum_{i=0}^{e-1}X_iZ^i$ on $F[Z]/(g)$.

For a nonzero tuple $\mathbf x\in F^e$, put
$r=\min_i v_{\mathfrak p}(x_i)$. After dividing its coordinates
by $\pi^r$, their reductions are not all zero. Multiplication by
their reduced polynomial is invertible in the residue field
extension defined by $g$. Its determinant is therefore a unit.
By homogeneity,
\[
                 v_{\mathfrak p}(N_g(\mathbf x))=er.
\]
Consequently
\[
 A_e(\mathbf X_0,\ldots,\mathbf X_{e-1})
       =\sum_{j=0}^{e-1}\pi^jN_g(\mathbf X_j)
\]
is anisotropic: the nonzero summands have pairwise distinct
valuations modulo $e$, so the summand of least valuation cannot
cancel. This is a homogeneous degree-$e$ form in $e^2$ coordinates.
\end{proof}

\begin{theorem}\label{fa:number-field}
Replace $\mathbb Z$ by $\mathcal O_F$ and $\mathbb Q$ by a
number field $F$ in the finite algebra presentation of
Section~\ref{sec:explicit-finite-algebra}. Retain the same rank $n$,
guard $G$ and $r$ closed residuals. There is an explicitly specified
polynomial in the same base coordinates and one additional coordinate
whose $F$-solutions have exactly the required input projection.
\end{theorem}
\begin{proof}
All multiplication matrices, determinants, resultants, reciprocal
polynomials and trace residuals in
Theorem~\ref{thm:explicit-guarded-finite-algebra} are defined over
$\mathcal O_F$ by the same formulas. For fixed $\ell,d$, replace
the sum of squares defining $Q_{\ell,d}$ by an anisotropic form
applied to the list
\[
 \Psi(U),\quad a_{d+1},\ldots,a_n,\quad
 \widetilde R_1(U),\ldots,\widetilde R_r(U).
\]
The list has $1+n-d+r\le n+r$ entries. Take
$e=\lceil\sqrt{n+r}\rceil$, use the form of
Lemma~\ref{fa:anisotropic}, and pad the list with zeros.
Its vanishing is equivalent to the vanishing of every listed entry.

The proof of the finite algebra theorem now applies verbatim:
it uses characteristic zero, the finite bound on bad separating
integers, and the one-dimensional kernel at a simple rational root.
The sole use of the ordering of $\mathbb Q$ was to combine the
listed residuals, and that step has been replaced by anisotropy.
Multiplication over all the finite choices gives the required
polynomial. Its coefficients belong to $\mathcal O_F$.
\end{proof}

For each factor, if the largest degree of the listed residuals is
$D_{\ell,d}$, its replacement has degree at most
$\lceil\sqrt{n+r}\rceil D_{\ell,d}$. Over a number field with
a real embedding, sums of squares give instead the bound
$2D_{\ell,d}$. Both methods preserve the number of coordinates.
The combining form depends on $F$; this
theorem does not assert that the original rational polynomial has
the same projection after arbitrary base change.
\end{document}
